\chapter{Introdução}
\label{cap:introducao}

Dentre as diversas atividades econômicas no semiárido brasileiro, a apicultura se destaca por apresentar inúmeras vantagens. As condições climáticas dessa região, aliadas ao potencial da vegetação nativa, influenciam na produção de méis com características únicas e alta qualidade, sendo bastante valorizados. Essa prática também se destaca por seu caráter essencialmente ecológico, unindo sustentabilidade e rentabilidade, pois contribui para a manutenção da biodiversidade por meio da polinização. Além disso, trata-se de uma atividade que exige baixo investimento e custo de implementação, sendo facilmente integrada como fonte complementar de renda nas comunidades rurais \cite{veras2025}. Sendo assim, a apicultura se consolida como uma importante fonte de renda, oferecendo maior segurança financeira ao homem do campo em um contexto em que as atividades agrícolas são frequentemente incertas \cite{araujo2024}.

Entretanto, apesar do cenário ambiental e econômico favorável, a plena conversão desse potencial em lucro real enfrenta barreiras estruturais. Por ter um forte caráter familiar, muitos apicultores acabam tendo dificuldades para competir no mercado, mesmo quando produzem um mel de alta qualidade e bastante valorizado \cite{veras2025}.

Nesse sentido, o cooperativismo e o associativismo surgem como alternativas que ajudam a diminuir custos e a melhorar o acesso a recursos tecnológicos, aumentando a competitividade dos produtores \cite{barbosa2020}. Essa forma de organização se relaciona com os princípios da Economia Solidária, que propõe o trabalho baseado na associação entre iguais, com foco na cooperação e na igualdade entre os membros, fortalecendo o grupo em relação ao indivíduo. Nesse modelo, o sucesso do empreendimento está na autogestão e na participação democrática, tornando a transparência e a justiça na divisão dos resultados pilares essenciais para manter a confiança e a coesão do grupo \cite{singer2002}.

Contudo, a transição desse ideal solidário para a prática apresenta desafios significativos. Conforme aponta \citeonline{singer2002}, a manutenção de um poder de decisão estritamente igualitário para todas as questões pode impactar a agilidade dos processos internos, uma vez que demandas de menor escala, que exigem respostas rápidas, acabam levando mais tempo do que o necessário. Assim, adota-se como solução a criação de uma diretoria democrática, mantendo-se alinhada à idealização solidária, mas com autonomia para deliberar sobre questões simples, evitando gargalos e tornando os processos mais ágeis \cite{aguiar2026}.

Portanto, embora essa centralização operacional de funções contribua significativamente para a gestão em associações apícolas, ela gera uma inevitável assimetria informacional. Como nem todos os membros participam ativamente da rotina administrativa diária, cria-se um distanciamento natural no acesso direto às informações sobre o que ocorre dentro do empreendimento. Esse cenário, se não for mediado por uma comunicação eficiente, tende a gerar desconfiança no grupo, fragilizando a autogestão e, em última instância, a própria sustentação da associação, uma vez que a transparência é o pilar que sustenta a união entre os associados \cite{bertolin2008}.

Para mitigar esses riscos e preservar a coesão do grupo, é essencial que a associação adote mecanismos de transparência total em suas ações e no relato de seus resultados. Assim, a tecnologia surge como uma garantia de que os princípios da Economia Solidária sejam respeitados, assegurando que a informação flua de maneira justa e que todos os membros tenham segurança sobre a integridade do empreendimento coletivo \cite{aguiar2026,aprova2025}.

Todavia, a adoção de uma ferramenta tecnológica não assegura, por si só, a resolução desses impasses nas associações. Conforme destaca Ian Sommerville, muitos sistemas falham por não atenderem às necessidades dos usuários ou por não serem confiáveis. Quando o software apresenta erros, inconsistências ou é de difícil utilização, tende a ser rejeitado pelos associados, pois compromete a confiança nas informações apresentadas \cite{sommerville2019}.

Dessa forma, para que a tecnologia sustente a autogestão, torna-se fundamental que o sistema seja desenvolvido com processos bem definidos, priorizando a confiabilidade e a facilidade de manutenção, de modo que os resultados sejam percebidos como fiéis à realidade \cite{iso25010}.

Nesse sentido, a Engenharia de Software fornece princípios, processos e práticas que auxiliam no desenvolvimento de aplicações mais consistente, reduzindo a ocorrência de falhas e aumentando a qualidade do produto entregue \cite{sommerville2019}. Entre essas práticas, destaca-se o Desenvolvimento Orientado por Testes (Test-Driven Development), abordagem que incentiva a especificação e a validação das regras de negócio por meio da criação de testes antes da implementação das funcionalidades. Com isso, busca-se assegurar que as funcionalidades atendam aos requisitos estabelecidos e mantenham seu comportamento esperado ao longo da evolução do sistema \cite{beck2002}.

\section{Objetivos}
\label{sec:objetivos}

\subsection{Objetivo Geral}

O objetivo geral deste trabalho consiste em desenvolver uma aplicação web voltada para a automação das rotinas internas da Associação de Apicultores de Poranga, situada no semiárido cearense. Busca-se, por meio de práticas da Engenharia de Software, construir uma ferramenta confiável que preserve a integridade das informações, confira clareza à gestão administrativa e mitigue a assimetria informacional, apoiando diretamente o princípio da autogestão e a sustentabilidade da associação.

\subsection{Objetivos Específicos}

Para alcançar o objetivo geral proposto, foram estabelecidos os seguintes objetivos específicos:

\begin{itemize}
    \item \textbf{Levantar e especificar os requisitos} funcionais e não funcionais junto aos membros e à diretoria da Associação de Apicultores de Poranga, mapeando os gargalos das rotinas manuais atuais;
    \item \textbf{Projetar a arquitetura do sistema web}, definindo a modelagem do banco de dados e a estrutura da aplicação de modo a garantir manutenibilidade e escalabilidade;
    \item \textbf{Implementar a aplicação web} utilizando práticas ágeis e padrões de projeto que priorizem a clareza do código e a separação de responsabilidades;
    \item \textbf{Aplicar o Desenvolvimento Orientado por Testes (TDD)}, construindo uma suíte de testes automatizados para validar as regras de negócio antes e durante a implementação, assegurando a confiabilidade do software;
    \item \textbf{Avaliar a qualidade do sistema} com base em critérios de usabilidade e integridade de dados, garantindo que a ferramenta seja facilmente adotada pelos associados e reflita com fidelidade a realidade administrativa da instituição.
\end{itemize}